The dataset and topic-selection procedures were designed to reflect the intended deployment conditions for graph-aware prompting while respecting confidentiality constraints on private materials. Concretely, our experimental materials comprised two complementary source types: (i) internal research notes and short-form documentation maintained by collaborating teams, and (ii) a small, curated collection of PDF documents intended to serve as external evidence. Details that could identify individual projects or reveal proprietary information are withheld; below we describe the selection criteria, curation workflow, and the unitization strategy used for retrieval and generation so that the experimental protocol can be evaluated and (where permitted) reproduced.

Source inclusion and exclusion. Documents drawn from internal notes were included only after explicit access approval from the originating teams and a light relevance screening against the target topics. The curated PDF collection was assembled to provide topical breadth and to include exemplar evidence artifacts (e.g., short papers, technical reports). Exclusion criteria removed documents that were manifestly out of scope (unrelated domains), near-duplicates, or documents with licensing restrictions that prevented use in experiments. We note that provenance tracking, archiving, and data citation are widely regarded as important features for curated databases and similar collections; issues of provenance and resource constraints in small curated projects have been documented in prior work on curated scientific databases \cite{web_fowler_2020_cross_tier_web_programming_for_curated_databases}.

Evidence curation workflow. Each candidate source underwent a lightweight curation pipeline comprising metadata normalization (title, authors, year, provenance label), text extraction (OCR or PDF parsing where necessary), and quality checks for extraction errors. Extracted text was segmented into retrievable units and indexed using the same embedding and ranking pipeline applied at retrieval time. When a source contained multiple distinct sections or appendices, these components were preserved as separate indexed items to avoid conflating disparate evidence. For private-note items, curators additionally recorded access and usage approvals and applied redaction where required by agreements. Because of privacy constraints, raw curated artifacts are not publicly released with this paper; the protocol above captures the operational steps taken.

Retrieval-unit design and evidence curation. Following best practices in hierarchical evidence assembly and curation, we indexed document fragments at section- and paragraph-level granularity so that retrieval could return focused passages rather than whole documents; related approaches emphasize hierarchical retrieval and post-retrieval evidence curation to remove irrelevant or near-duplicate passages \cite{web_choe_2025_hierarchical_retrieval_with_evidence_curation_fo}. Candidate units were annotated with provenance metadata and assigned stable identifiers; these identifiers were referenced by the section-generation module when assertions required supporting evidence.

Topic selection rationale. Two topical domains were chosen to exercise complementary challenges for retrieval-guided generation. One domain emphasized relatively dense, well-structured evidence (higher signal-to-noise in retrieved passages), while the other emphasized sparser or more heterogeneous evidence sources (higher retrieval difficulty). This contrast was intended to probe how graph-guided sectioning performs under varying evidence availability. (General background knowledge: selecting domains that vary in evidence density helps surface robustness differences between retrieval-guided generation methods.)

Definition of document units for retrieval and generation. Consistent with the document-graph formulation (Section 3.1), the basic unit for retrieval and generation was a section node. Each node carried metadata including a unique identifier, heading text, a target length estimate, and a role label (for example: Background, Method, Result, Discussion). During preprocessing, source texts were segmented into candidate evidence units aligned to typical section- and paragraph-level boundaries; these candidate units were the items indexed and returned by the retrieval step for a given node's queries. Generators were required to produce JSON-conformant section outputs that reference evidence-unit identifiers when asserting facts or claims, enabling downstream auditing and automatic checks.

Mapping units to queries and conservative fallbacks. For each section node the retrieval query was constructed from a compact contextual representation: the node heading, a short abstract of preceding nodes (to supply local context), explicit claim templates where applicable, and a configurable set of role-derived keywords. Retrieved items were ranked and filtered before being injected into the prompt for constrained generation. When the retrieval hit set was empty or judged of low quality by automated heuristics, the generation prompt included explicit conservative-language instructions and required that any unverifiable claim be flagged as such in the JSON output.

Privacy, access, and provenance disclosures. Access to internal notes was limited to approved experimenters and governed by data-use agreements; accordingly, the internal-note corpus is not publicly released. The curated PDF collection contains items that are either public-domain or licensed for redistribution within the project; where redistribution is blocked, we provide provenance metadata and examples of the extraction and indexing artifacts produced during curation. All human annotation and reviewer-evaluation procedures involving private materials followed the consent and confidentiality procedures described in Section 4.6.

Limitations. Because parts of the dataset are private and the curated collection is intentionally small-scale, the experimental results should be interpreted with caution regarding broad generalization. The curation and topic-selection process prioritized realism for the intended application (writing support for domain teams) over large-scale public benchmarking; we discuss these limitations and potential mitigations (e.g., open benchmark construction and larger curated collections) in Section 6.